\section{Related work}
\label{sec:related_work}

\textbf{Preventing Eavesdropping with Microphone's Nonlinearity. }
Several existing works attempt to jam microphones using ultrasound. Roy~et al. inject white noise to microphones using ultrasound \cite{backdoor}. Li~et al. generate noise with variable frequency according to a preset key to prevent the unauthorized devices from recording audios, while enabling the authorized users to recover speech signals from the noisy recordings \cite{patronus}. Sun~et al. propose MicShield, which can prevent the always-on microphones in smart home devices from recording private speech while passing the preset voice commands \cite{sun_alexa_2020}. Chen~et al. integrate ultrasonic transmitters into a wearable bracelet to expand the effective jamming coverage \cite{CHI}. However, the noises used in these works are either single tones with variable frequency or white noise, which are not robust when facing speech enhancement methods as demonstrated by the experiments in this paper.


\textbf{Other Applications with Microphone's Nonlinearity. }
Many works explore the microphone nonlinearity for other purposes, such as inaudible voice commands injection \cite{10.1145/3133956.3134052, 211283, yan_surfingattack_2020}, defense of inaudible commands \cite{he_canceling_2019, Zhang2021EarArrayDA, 211283}, communication \cite{backdoor}, and authentication \cite{zhou_nauth_2019}. Yan~et al. transmit the inaudible voice commands through solid medium to improve the stealthiness of the attack \cite{yan_surfingattack_2020}. Roy~et al. expand the attack range by striping different frequency bands to different transmitters \cite{211283}. Zhang~et al. exploit the difference of propagation attenuation between ultrasound and voice to distinguish inaudible commands injection from normal voice commands. Zhou~et al. validate that the parameters of the nonlinearity model in different microphones can be used as features for device authentication \cite{zhou_nauth_2019}.